% ====================================================================
%  Chapter 1 — Introduction
% --------------------------------------------------------------------
%  Template chapter. Replace every bracketed [LIKE_THIS] placeholder
%  with the actual value before submission. The chapter sections follow
%  the standard MPhil/journal-paper introduction structure:
%    1.1 Motivation
%    1.2 Problem statement
%    1.3 Research questions
%    1.4 Contributions
%    1.5 Thesis structure
%  Citations resolve against the .bib block in RUN_GUIDE.md §16.
% ====================================================================

\chapter{Introduction}
\label{ch:introduction}

\section{Motivation}
\label{sec:motivation}

Speaker verification (SV) — the task of deciding whether two audio
recordings come from the same speaker — is a foundational biometric
technology underpinning voice authentication, forensic audio analysis,
and conversational-AI personalisation. The dominant publicly available
benchmarks (VoxCeleb1 \cite{nagrani2017voxceleb}, VoxCeleb2, NIST SRE
series) are constructed almost exclusively from English speech. As a
result, the state-of-the-art has converged on architectures and
training recipes that, while highly effective on English, are not
guaranteed to generalise to low-resource languages.

Sri Lanka presents a particularly underserved case. The country's two
primary languages —~Sinhala (an Indo-Aryan language with roughly
17 million speakers) and Tamil (a Dravidian language with roughly
5 million Sri Lankan speakers) —~have, prior to this work and our
group's SLCeleb corpus \cite{slceleb2025}, no publicly catalogued SV
benchmarks and no published SV baselines. Existing
multilingual self-supervised speech encoders cover them only as part
of broader low-resource batches \cite{chen2022wavlm,
baevski2020wav2vec2}, with no targeted SV evaluation. The Sri Lankan
deployment context further compounds the difficulty: bilingual and
code-switched speech is normal in everyday usage, telephony channels
dominate practical deployment, and labelled speakers are expensive to
collect.

\section{Problem statement}
\label{sec:problem}

This thesis addresses the question of how to construct a speaker
verification system that performs reliably on Sinhala, Tamil, and
code-switched audio when only a small labelled corpus is available.
We frame the problem along three dimensions:

\begin{enumerate}
    \item \textbf{Architectural choice.} Modern SV architectures
    (ECAPA-TDNN \cite{desplanques2020ecapa}, ResNet-SE variants,
    self-supervised front-ends) were validated on English VoxCeleb
    scale. Whether their inductive biases hold for low-resource
    Sinhala/Tamil is an open empirical question.
    \item \textbf{Cross-lingual transfer.} When a target-language
    corpus is small ($\leq$~10\,000 utterances), training from
    scratch overfits. Cross-lingual fine-tuning from an English-
    pretrained checkpoint is the natural alternative; the right
    freezing, learning-rate, and regularisation policy for SV
    fine-tuning specifically remains under-studied.
    \item \textbf{Calibration and operating point.} Per-language
    decision thresholds, score normalisation (AS-Norm
    \cite{matejka2017analysis}), generative scoring backends (PLDA),
    and language-aware regularisation (domain adversarial training
    \cite{ganin2015dann}) are standard tools for cross-lingual SV in
    the literature, but their interactions on Sinhala/Tamil data
    are not characterised.
\end{enumerate}

\section{Research questions}
\label{sec:rqs}

We address the following research questions:

\begin{description}
    \item[\textbf{RQ1}] How well do released, English/VoxCeleb-trained
    SV models transfer \emph{zero-shot} to Sinhala and Tamil, and how
    large is the degradation relative to their published English
    benchmarks?
    \item[\textbf{RQ2}] Which fine-tuning policy — frozen-encoder
    probing, parameter-efficient adaptation (LoRA), full fine-tuning,
    or layer-wise LR decay — best adapts a pretrained speech
    foundation model to Sinhala/Tamil per trainable parameter?
    \item[\textbf{RQ3}] Does language-adversarial training (gradient
    reversal / DANN \cite{ganin2015dann}) improve cross-lingual
    verification of bilingual speakers, and by how much?
    \item[\textbf{RQ4}] At eval time, do score normalisation (AS-Norm)
    and a generative scoring backend (PLDA) compose constructively, or
    do they redundantly model the same residual variance?
    \item[\textbf{RQ5}] How large is the gap between monolingual
    (Sinhala-only, Tamil-only) and cross-lingual EER for a single
    trained model, and can it be closed via threshold calibration
    alone or does it require architectural intervention?
\end{description}

\section{Contributions}
\label{sec:contributions}

The contributions of this thesis are:

\begin{enumerate}
    \item \textbf{A reproducible Sinhala/Tamil SV benchmark.}
    Benchmark v0 covers 527 speakers (478 Sinhala from OpenSLR SLR52,
    49 Tamil from OpenSLR SLR65) with seed-reproducible monolingual
    trial lists (12{,}000 trials per language, gender-controlled
    impostors where metadata permits), produced by a released
    deterministic pipeline; benchmark v1 extends the evaluation to
    SLCeleb \cite{slceleb2025}, our group's 280-speaker in-the-wild
    corpus with multi-session and bilingual speakers.
    \item \textbf{A comprehensive empirical study} of cross-lingual
    transfer for SV under low-resource constraints, in three suites:
    the first zero-shot baseline table for Sinhala (four modern
    architectures); a training-free adaptation ablation (AS-Norm,
    PLDA, per-language calibration) that quantifies how much of the
    cross-lingual degradation is a fixable score-domain problem; and
    a four-policy fine-tuning comparison (frozen / LoRA / full /
    layer-wise LR decay) of a WavLM foundation model, replicated
    over three seeds.
    \item \textbf{Per-language and pooled evaluation infrastructure}
    that supports principled comparison of Sinhala-only,
    Tamil-only, and cross-lingual trial sets in a single trainer
    invocation, with paper-grade statistical reporting
    (mean$\pm$std over three random seeds).
    \item \textbf{A practical open-source toolkit} extending the
    \texttt{voxceleb\_trainer} codebase \cite{slspv2026repo} with
    eleven new features (FEATURE-001..FEATURE-011) and twenty-six
    bug fixes (BUGFIX-001..BUGFIX-026). Each feature ships with a
    design document and is selectable via a single YAML key, enabling
    independent ablation.
\end{enumerate}

\section{Thesis structure}
\label{sec:structure}

The remainder of the thesis is organised as follows.
Chapter~\ref{ch:related} reviews prior work on speaker verification
architectures, cross-lingual transfer, score normalisation, and
language-aware modelling.
Chapter~\ref{ch:methods} describes the corpus, architectures, training
recipe, scoring backends, and evaluation protocol.
Chapter~\ref{ch:results} reports the headline and ablation results.
% Chapters 5 (Discussion) and 6 (Conclusion) are not scaffolded yet;
% switch the plain numbers below back to \ref{ch:discussion} /
% \ref{ch:conclusion} once those files exist.
Chapter~5 analyses the findings, addresses the
research questions of Section~\ref{sec:rqs}, and discusses limitations.
Chapter~6 summarises the contributions and outlines
future work, including the SLCeleb (v1) open-set benchmark, the
bilingual cross-lingual study, and self-supervised pretraining on
unlabelled Sri Lankan audio.
